\section{Related Work}
\label{sec:related}

Three converging research threads form the prior art for \sysname; we
summarise each with its gap, then position \sysname against all three.

\smallskip
\noindent\textbf{Consent management on distributed ledgers.}
MedRec~\cite{azaria2016medrec} records patient consent immutably
on-chain---the durability and owner-exclusive control we inherit---but
treats the record as the protocol \emph{endpoint}: no training is gated
by it, and nothing shows a withdrawal produced a model update
(semantic-web vocabularies like GConsent formalise consent off-chain
without enforcement). \emph{Gap}: consent is recorded but does not
influence training.

\smallskip
\noindent\textbf{Machine unlearning.}
SISA~\cite{bourtoule2021sisa} pre-shards training data so that removal
retrains only the affected shard; FedEraser~\cite{liu2021federaser}
replays federated aggregation history with a calibration step; and
gradient-ascent methods take a few negative-learning-rate steps on the
data to be removed. All are validated on standalone benchmarks, but
none specifies how unlearning is triggered, by whom, or how a deployer
demonstrates it actually happened. \emph{Gap}: unlearning is
implementable but not verifiable end-to-end.

\smallskip
\noindent\textbf{Blockchain-based federated learning.}
BlockFL~\cite{kim2020blockchained} logs device-side updates on-chain
against malicious aggregators and BFLC~\cite{li2020bflc} replaces the
aggregator with committee consensus, but neither encodes consent: a
joined participant cannot withdraw without trusting the orchestrator.
Membership inference~\cite{shokri2017mia,carlini2022lira,choquettechoo2021labelonly}
is the canonical test for residual memorisation, yet no blockchain-FL
system uses MIA scores as an on-chain acceptance criterion. \emph{Gap}:
updates are auditable but consent semantics and forgetting verification
are absent.

\smallskip
\noindent\textbf{Positioning.}
Across the four implicated dimensions---consent representation,
unlearning trigger and method, verification, and on-chain
footprint---each prior system covers at most one natively. \sysname is
the first design we are aware of to address all four in a single
chain-enforced protocol at gas-free cost: it adopts MedRec's registry
pattern (vertex~A) under the \clsm safety and liveness properties,
accepts any of the three unlearning paradigms as a black box
(vertex~B), and makes three-MIA verification (vertex~C) the chain's
gating condition for a consent transition. The resulting loop---consent
change $\rightarrow$ unlearning $\rightarrow$ verification
$\rightarrow$ attestation $\rightarrow$ finalisation---is the missing
edge.

The choice of \emph{which} blockchain matters at IoT scale:
permissionless gas-based chains (Ethereum, Polygon) impose recurring
per-operation costs that scale linearly with device count and consent
operation rate. Section~\ref{sec:evaluation} quantifies this against
PureChain~\cite{kim2025purechain}, a public, gas-free, EVM-compatible
chain with PoA$^2$ consensus and Merkle-proof-backed state. PureChain
is not the contribution of this paper; rather, we show the protocol is
operationally identical on both chain types while the economic model
is fundamentally different.
